\documentclass[11pt,a4paper]{article}
\usepackage[margin=2.5cm]{geometry}
\usepackage{amsmath,amssymb,physics}
\usepackage{booktabs}
\usepackage{longtable}
\usepackage{array}
\usepackage{hyperref}
\usepackage{enumitem}
\usepackage{tcolorbox}
\tcbuselibrary{skins,breakable}

\newtcolorbox{decisionbox}[1][]{enhanced,breakable,
  colback=green!5!white,colframe=green!50!black,
  fonttitle=\bfseries,title={#1},boxrule=0.6pt,arc=4pt}
\newtcolorbox{conflictbox}[1][]{enhanced,breakable,
  colback=red!5!white,colframe=red!50!black,
  fonttitle=\bfseries,title={#1},boxrule=0.6pt,arc=4pt}
\newtcolorbox{notebox}[1][]{enhanced,breakable,
  colback=blue!4!white,colframe=blue!40!black,
  fonttitle=\bfseries,title={#1},boxrule=0.6pt,arc=4pt}

\hypersetup{colorlinks=true,linkcolor=blue,citecolor=blue,urlcolor=blue}

\title{\textbf{SSV Units \& Length-Scale Audit}\\[0.4em]
  \large Inventory, Resolution, and Defined Reference Table}
\author{Working document for Papers I and II}
\date{\today}

\begin{document}
\maketitle

\begin{abstract}
This is a working document, not a paper.  Its purpose is to identify every
appearance of the symbols $\xi$, $m_0$, $a$, $R^*$, $\bar\lambda_e$,
$\xi_{\rm EW}$, and $\ell_P$ in the SSV manuscripts (Papers I and II), audit
each for consistency, and produce a single defined reference table that
subsequent revisions can use as the canonical source.  The motivation is a
concrete inconsistency identified during the draft-22 read-through of
Paper II: the symbol $\xi$ is used in at least three numerically incompatible
ways across the manuscripts.

Tracing this inconsistency leads to a substantive clarification of the
framework's ontology, captured in six decisions.  The single symbol $\xi$
has been doing the work of two distinct concepts: the \emph{electron
healing length} $\hbar/(m_e c)$, which sets the electron defect's spatial
extent, and the \emph{medium domain-of-validity scale} $\ell_{\rm grain}$,
below which the LogSE continuum description ceases to be self-consistent.
We separate them, and identify $\ell_{\rm grain}$ with the heaviest stable
defect mass --- the proton's Compton scale, $a_p \approx 2\times
10^{-16}\,\mathrm{m}$.  This recasts the framework as a continuum
effective theory in the same epistemic register as Navier--Stokes, with a
definite domain of validity and explicit silence about deeper microscopic
substrate.

The audit further clarifies three structural points.  First, the Higgs
is not a fundamental medium length scale but a denting event: the
125\,GeV LHC resonance is the energy required to drive a localised
$|\Psi|\to 0$ cavity of specific geometry against the saturation pressure
$P_0$.  Second, correlations within the medium's saturated coherent
state --- including those exhibited by entangled particles --- are
holographic in character and have no propagation speed; the
``speed-of-spooky-action'' lower bounds (Salart 2008, Yin 2013) constrain
hypothetical signal velocities and are vacuous in the present framework.
Third, all radiation between defects --- both electromagnetic and
gravitational --- propagates in the transverse channel at $c_\perp$,
because mass produces a density depression (a transverse deformation of
the saturated state) rather than a compression, and the longitudinal
channel is gapped at the Higgs amplitude scale.  The longitudinal speed
$c$ remains an internal medium parameter (setting saturation pressure,
the Higgs gap, and denting-event geometry) but is not the propagation
speed of any astrophysical signal.  This resolves the apparent GW170817
problem: gravitational and electromagnetic radiation share the carrier
channel and arrive simultaneously up to source-physics timescales.

The audit produces a canonical defined-symbol table to anchor manuscript
revisions, a clear domain-of-validity statement for SSV, and a list of
genuine open physics questions that the audit could not itself resolve.
\end{abstract}

\tableofcontents
\newpage

% ═══════════════════════════════════════════════════════════════════════
\section{Method}
% ═══════════════════════════════════════════════════════════════════════

We do four things, in this order:
\begin{enumerate}
\item \textbf{Inventory.} List every appearance of each symbol in Paper~I and
      Paper~II, with line number and context.  No interpretation at this stage.
\item \textbf{Read.} For each symbol, identify what value or relation each use
      assigns or implies.  Flag conflicts.
\item \textbf{Resolve in dependency order.}  Some symbols are defined in
      terms of others; we work outwards from the most fundamental.  The order
      is: $m_0$ first (since it sets the medium grain); then $\xi$ (which
      depends on $m_0$); then $a$ and $R^*$ (which depend on $\xi$); then
      $\xi_{\rm EW}$ (which compares against $\xi$); then $\ell_P$ (whose
      role in SSV must be decided).
\item \textbf{Tabulate.}  Produce a single defined-symbol table to be used
      as the canonical reference for revisions.
\end{enumerate}

We do \emph{not} attempt at this stage to fix the manuscripts.  The output
of this audit is the table at the end; the manuscript revisions follow once
the table is settled.

% ═══════════════════════════════════════════════════════════════════════
\section{Inventory}
% ═══════════════════════════════════════════════════════════════════════

Line numbers refer to \texttt{SSV-01-draft-07.tex} (Paper~I) and the working
draft of Paper~II.  Square brackets give an abridged context.

\subsection{Uses of $m_0$}

\begin{longtable}{@{}llp{9cm}@{}}
\toprule
Paper & Line & Context \\
\midrule
\endhead
I  & 48   & $b = m_0 c^2/(2\rho_0)$, $\xi = \hbar/(m_0 c)$ [stiffness fixing] \\
I  & 189  & $-(\hbar^2/2m_0)|\nabla\Psi|^2$ [LogSE kinetic term] \\
I  & 205  & $-(\hbar^2/2m_0)\nabla^2\Psi$ [LogSE] \\
I  & 212  & $\mathbf{v}=(\hbar/m_0)\nabla\theta$ [Madelung velocity] \\
I  & 222--223 & Madelung momentum equation, $\nabla P/m_0 \rho$ \\
I  & 231--232 & $c_{\rm thermo}=\sqrt{(1/m_0)dP/d\rho}=\sqrt{b/m_0}$;
                $c_s=\sqrt{2b\rho_0/m_0}$ \\
I  & 247  & $\xi = \hbar/\sqrt{2 m_0 b\rho_0} = \hbar/(m_0 c)$ \\
I  & 250  & ``the Compton wavelength of the reference mass $m_0$'' \\
I  & 268  & $c_\perp/c = \sqrt{\lambda m_0/\rho_0} \equiv \alpha$ \\
I  & 274--277 & $\lambda = \alpha^2 \rho_0/m_0$;
                $\rho_0/m_0 \approx 3.75\times 10^7$ \\
I  & 322--323 & ``$m_0$ is a reference mass scale (self-consistently $m_e$
                for the electron)'' \\
I  & 1083  & $\kappa_0 = 2\pi\hbar/m_e$, $\xi = \hbar/(m_e c)$,
                $R_e = \hbar/(\alpha m_e c)$ [explicit identification
                $m_0 \to m_e$] \\
\midrule
II & 158  & ``$\xi = \hbar/(m_0 c)$'' [Framework recap] \\
II & 304  & $\partial_t^2\delta\rho = c^2\nabla^2\delta\rho - (\hbar^2/4m_0^2)\nabla^4\delta\rho$
            [Bogoliubov dispersion] \\
II & 765  & $E_{\rm recon} \sim \hbar_{\rm eff}^2 \rho_0/(m_0^2 a^2)$
            [W mass scalar estimate] \\
II & 768  & $a \sim \hbar_{\rm eff}/(m_e c)$ [identification: vortex core
            scale set by $m_e$] \\
II & 1645 & $g^{\rm eff}_{\mu\nu} = (\rho_0/(m_0 c)) \cdot \dots$
            [acoustic metric] \\
\bottomrule
\end{longtable}

\subsection{Uses of $\xi$}

\begin{longtable}{@{}llp{9cm}@{}}
\toprule
Paper & Line & Context \\
\midrule
\endhead
I  & 48   & $\xi = \hbar/(m_0 c)$ [framework definition] \\
I  & 50   & $R_e^* = \xi/\alpha$ [equilibrium ring radius] \\
I  & 247  & $\xi = \hbar/(m_0 c)$ [explicit derivation] \\
I  & 354  & $a^* = \xi$ [vortex core size from minimisation] \\
I  & 395  & $R_e^* = \xi/\alpha = \hbar/(\alpha m_0 c)$ [boxed result] \\
I  & 475--478 & $\xi \sim \hbar/(m_e c)$ [classical electron radius
                $r_e = \alpha\xi$] \\
I  & 1083 & $\xi = \hbar/(m_e c)$ [explicit numerical value] \\
\midrule
II & 158  & $\xi = \hbar/(m_0 c)$ [framework recap, line~161 of working] \\
II & 162  & ``$\alpha = c_\perp^2/c^2 \approx 1/137$'' [no $\xi$ statement
            but adjacent] \\
II & 402  & $R^* = \xi/\alpha$ [equilibrium ring radius] \\
II & 415  & $\xi \sim \ell_P$ [\textbf{conflict: $\xi$ identified with
            Planck length}] \\
II & 845--846 & $\xi = \hbar/(m_e c) \approx 3.86\times 10^{-13}\,\mathrm{m}$
                [Higgs hierarchy: vortex core healing length] \\
II & 900  & $a \sim \xi/\alpha$ [reconnection scale; \textbf{$a$ here is
            ring-scale, not core-scale}] \\
II & 910  & ``minor radius reaches $\xi$'' [Rayleigh jet pinch-off; $\xi$ is
            here a sub-Compton scale] \\
II & 1491 & $\xi = \hbar/(m_e c)$, $a_p/\xi = m_e/m_p$
            [proton sub-grain ratio] \\
II & 1612 & $\xi = \hbar/(m_e c) \approx 3.86\times 10^{-13}\,\mathrm{m}$
            [graviton-free, UV cutoff argument] \\
\bottomrule
\end{longtable}

\subsection{Uses of $a$ (vortex core radius)}

\begin{longtable}{@{}llp{9cm}@{}}
\toprule
Paper & Line & Context \\
\midrule
\endhead
I  & 50   & ``core radius $a$'' [generic] \\
I  & 354  & $a^* = \xi$ [minimisation result; \textbf{core radius equals
            healing length}] \\
\midrule
II & 176  & ``$a \ll R_e$'' [thin-core assumption for far field of
            electron] \\
II & 660--668 & $\mathcal{E}_\ell = \tfrac{1}{2}\rho_0\kappa_0^2 \ln(L/a)$;
                $L \sim 1$\,fm [string tension, $a$ = strand core radius] \\
II & 676  & $r_c \sim a/\alpha$ [confinement scale] \\
II & 692  & ``$r \lesssim a$'' [asymptotic freedom, vortex core] \\
II & 749--750 & ``$r \lesssim a$ — hence the short range of the weak
                force'' [reconnection requires core overlap] \\
II & 765  & $a$ in $E_{\rm recon}$ [W mass scalar, identified with core] \\
II & 768  & $a \sim \hbar_{\rm eff}/(m_e c)$
            [\textbf{core $a$ identified with $\bar\lambda_e$, not with
            $\xi/\alpha$}] \\
II & 900  & ``two vortex defects approach within one core radius $a \sim
            \xi/\alpha$'' [\textbf{contradicts II:768; here $a$ is
            ring-scale not core-scale}] \\
\bottomrule
\end{longtable}

\subsection{Uses of $R^*$, $R_e$, $\bar\lambda_e$}

\begin{longtable}{@{}llp{9cm}@{}}
\toprule
Paper & Line & Context \\
\midrule
\endhead
I  & 50, 395, 427, 1083, 1141, 1155 & $R_e^* = \xi/\alpha$
                  [electron equilibrium ring radius, recurring] \\
I  & 822  & $2R_e = 2\xi/\alpha$ [pion as flux-tube link spanning electron] \\
\midrule
II & 402  & $R^* = \xi/\alpha$ [g--2 finite-ring correction] \\
II & 419  & $\bar\lambda_e = \hbar/(m_e c)$ [reduced electron Compton] \\
II & 952  & ``the electron ring radius $R^* = \xi/\alpha$'' [Rayleigh jet] \\
II & 967  & ``$R_{\rm cap} = \phi R^*$'' [W cap geometry] \\
II & 1281 & ``$R^* = \xi/\alpha$'' [muon ring-breathing recap] \\
\bottomrule
\end{longtable}

\subsection{Uses of $\xi_{\rm EW}$ and $\ell_P$}

\begin{longtable}{@{}llp{9cm}@{}}
\toprule
Paper & Line & Context \\
\midrule
\endhead
II & 415  & $\ell_P$ as substitute for $\xi$ [\textbf{the central conflict}] \\
II & 420  & ``$\xi \sim \ell_P$ (the Planck length)'' [\textbf{explicit but
            inconsistent identification}] \\
II & 842  & $\xi_{\rm EW} = \hbar/(m_H c) \approx 1.6\times 10^{-18}\,$m
            [Higgs grain length] \\
II & 845--846 & $\xi_{\rm EW}/\xi \sim 10^{-5}$
                [hierarchy puzzle: with $\xi = \hbar/(m_e c)$] \\
II & 1612 & $\xi$ as UV cutoff $\approx 3.86\times 10^{-13}\,$m
            [\textbf{conflicts with $\xi \sim \ell_P$ at II:415}] \\
\bottomrule
\end{longtable}

% ═══════════════════════════════════════════════════════════════════════
\section{Reading the Inventory}
% ═══════════════════════════════════════════════════════════════════════

\subsection{What is consistent}

\textbf{Paper I is internally consistent} on the relation
$\xi = \hbar/(m_0 c)$ throughout, except at line~1083 where $m_0$ is
silently replaced by $m_e$.  That replacement is licensed by the framework
declaration at line~322--323 (``$m_0$ is a reference mass scale, self-consistently
$m_e$ for the electron'').  So Paper~I's working numerical convention is:
\begin{equation*}
    \xi^{\rm (I)}_{\rm working}
      = \frac{\hbar}{m_e c}
      \approx 3.86\times 10^{-13}\,\mathrm{m},
\end{equation*}
with the understanding that this is the value of $\xi$ \emph{evaluated at
$m_0 = m_e$}, appropriate for the electron defect.

\textbf{Paper II's force-sector and gravity content is consistent with this
working value}: the proton sub-grain ratio $a_p/\xi = m_e/m_p$
(line~1494), the Higgs hierarchy comparison (line~845), and the UV cutoff
argument (line~1612) all use $\xi \approx 3.86\times 10^{-13}\,\mathrm{m}$.

\subsection{What conflicts}

Two distinct conflicts.

\begin{conflictbox}[Conflict 1: $\xi$ vs.\ $\ell_P$ at line II:415]
The g--2 finite-ring correction (\S3.5 of Paper~II) substitutes
$\xi \to \ell_P$ to obtain a $\sim 10^{-44}$ suppression.  The Planck length
is $\ell_P \approx 1.6\times 10^{-35}\,\mathrm{m}$; the working value of $\xi$
elsewhere in the paper is $3.86\times 10^{-13}\,\mathrm{m}$.  The two differ
by 22 orders of magnitude.  This substitution is incompatible with every
other appearance of $\xi$ in Paper~II and Paper~I.  It appears to be a
fossil from an earlier convention in which $\xi$ was identified with the
Planck length.  The substitution must either be removed, or the rest of
the paper rewritten to use $\xi \sim \ell_P$ throughout (which would
invalidate the G consistency check, the proton sub-grain calculation, the
Higgs hierarchy comparison, and the UV-cutoff argument).
\end{conflictbox}

\begin{conflictbox}[Conflict 2: $a$ as core vs.\ ring scale]
The symbol $a$ is used in Paper~II to mean two different things:
\begin{itemize}
    \item In the strong-force section (\S4) and the W-mass scalar estimate
          (\S5.2, line~768), $a \sim \hbar/(m_e c)$ is the \emph{vortex core}
          radius, equal to $\xi$ at the working numerical value.
    \item In the Rayleigh-jet section (\S5.4, line~900), ``two vortex defects
          approach within one core radius $a \sim \xi/\alpha$'' uses $a$ as
          the \emph{ring} scale, equal to $R^*$.
\end{itemize}
These differ by a factor of $\alpha^{-1} \approx 137$.  In Paper~I (line~354)
the minimisation result is $a^* = \xi$, fixing the core radius equal to the
healing length.  Paper~II line~900 contradicts this.
\end{conflictbox}

\subsection{What is under-specified}

The framework's deepest assumption --- ``$m_0$ is a reference mass scale,
self-consistently $m_e$ for the electron'' (Paper~I, line~322--323) --- is
declared but never developed.  In particular:

\begin{itemize}
\item For the proton breather, what is $m_0$?  Paper~II implicitly takes
      $m_0 = m_e$ everywhere (so $\xi$ is always the electron Compton
      wavelength), but the proton has $m_p \approx 1836\, m_e$, and the
      proton's own ``Compton wavelength of the reference mass'' would
      naturally be $\hbar/(m_p c) = a_p$.  The relation $a_p/\xi = m_e/m_p$
      is therefore consistent with reading $\xi$ as the electron's $m_0$
      value and $a_p$ as the proton's $m_0$ value of the \emph{same}
      ``Compton wavelength of the reference mass'' formula.
\item For the medium itself, is there a third value of $m_0$ --- a true
      grain mass independent of any specific defect?  This is the
      ``Planck-equivalent'' intuition that motivated the original
      $\xi \sim \ell_P$ substitution at II:415.  The framework neither
      affirms nor denies it.
\end{itemize}

% ═══════════════════════════════════════════════════════════════════════
\section{Resolution in Dependency Order}
% ═══════════════════════════════════════════════════════════════════════

\subsection{Resolving $m_0$}

The framework's existing declaration (Paper~I, line~322--323) is correct
but incomplete.  $m_0$ is not a single universal constant of the medium;
it is a parameter that takes a specific value for each topological defect,
namely the rest mass of the lightest standing structure that defect
realises.  We make this explicit:

\begin{decisionbox}[Decision 1: $m_0$ is defect-relative]
The reference mass $m_0$ in any equation involving a single topological
defect is the rest mass of that defect's lightest standing realisation:
\begin{itemize}
\item Electron defect (toroidal vortex ring): $m_0 = m_e$.
\item Proton defect (knotted breather): $m_0 = m_p$.
\item Higgs amplitude mode (no defect; the canvas itself):
      $m_0 = m_H$ (with the caveat that the Higgs is not a localised
      structure but a global mode).
\end{itemize}
For derivations that span multiple defect types --- e.g.\ the proton
sub-grain coupling to the long-range phonon field, which involves both
$m_p$ (the proton breather) and a coarse-graining scale set by the
medium's coarsest stable defect ($m_e$) --- both reference masses appear
explicitly.  No single $m_0$ is universal.
\end{decisionbox}

This reading is consistent with every equation in Paper~I and Paper~II
where $m_0$ appears, provided the dual occurrences in II:768 ($a \sim
\hbar/(m_e c)$) and II:900 ($a \sim \xi/\alpha$) are read as different
quantities (see below).

\subsection{Resolving $\xi$: split into two symbols}

The simplest fix is to recognise that $\xi$ is doing the work of two
distinct quantities:

\begin{decisionbox}[Decision 2: $\xi$ is the electron healing length;
                    $\ell_{\rm grain}$ is the medium's domain-of-validity scale]
\begin{align}
    \xi              &\equiv \frac{\hbar}{m_e c}
                       \approx 3.86\times 10^{-13}\,\mathrm{m},
                       \tag{electron healing length} \\
    \ell_{\rm grain} &\equiv \frac{\hbar}{m_{\rm grain} c}
                       \approx 2.10\times 10^{-16}\,\mathrm{m},
                       \tag{medium domain-of-validity scale}
\end{align}
where $m_{\rm grain} = m_p$ is the heaviest stable defect mass (see
Decision~5 below).  The healing length $\xi$ is the SSV analogue of the
de~Broglie wavelength of the lightest stable defect; it sets the vortex
core size for electrons, the proton sub-grain coupling factor, and the
natural cutoff for atomic-scale calculations.  It is \emph{not} the
medium's UV cutoff in the deep sense.

The grain length $\ell_{\rm grain}$ is the scale at which the LogSE
continuum description ceases to be self-consistent.  Below this scale,
the medium supports only transient events (such as the Higgs and $W$
denting events of Paper~II \S5); no new stable defects exist.  The
framework is silent about what is ``below'' $\ell_{\rm grain}$ in the
same way Navier--Stokes is silent about what is below the molecular
scale of water --- this is a feature, not a defect, of an effective
continuum theory.  See \S\ref{sec:domain} below for the discussion of
domain of validity.
\end{decisionbox}

This decision implies the following revisions:
\begin{itemize}
\item Paper~II line~415 (the g--2 finite-ring correction): the substitution
      $\xi \to \ell_P$ is wrong.  The correct expansion variable is either
      $(\bar\lambda_e/R^*)^2 = \alpha^2$ (if the form factor cuts the loop
      integral off below the Compton scale) or a different small parameter
      that needs deriving.  This is a content-level fix, not a notation
      fix.  See \S\ref{sec:g2} below.
\item Paper~II \S6.6 (graviton-free / UV-cutoff argument): the cutoff is
      properly $\ell_{\rm grain}$, not $\xi$.  The current text claims the
      UV problem dissolves because momentum modes with $k \gg 1/\xi$ do
      not exist, but $1/\xi$ corresponds to the \emph{electron Compton
      momentum} $m_e c/\hbar$, well below QCD scales let alone Planck.
      The argument needs to be made about $\ell_{\rm grain}$ instead.
      With $\ell_{\rm grain} \approx 2\times 10^{-16}\,\mathrm{m}$, the
      cutoff is at the proton scale --- well below the Compton scale of
      the electron and far below the Planck scale.  See \S\ref{sec:domain}
      for the precise content of the argument.
\item Paper~II \S5.3 (Higgs hierarchy): the comparison
      $\xi_{\rm EW}/\xi$ is between the \emph{Higgs denting scale} and
      the electron Compton wavelength.  Since $\xi_{\rm EW}$ is itself
      not a structural medium scale but a derived event-energy length
      scale (see \S\ref{sec:higgs_event}), the ``hierarchy puzzle''
      framing in the existing text is misleading.  The genuine
      hierarchy question in SSV is the defect spectrum: why
      $m_e \ll m_p$, i.e.\ why is the toroidal vortex ring three orders
      of magnitude larger than the smallest stable knotted breather?
\end{itemize}

\subsection{Resolving $a$}

With Decision~1 in place, the apparent conflict at II:768 vs.\ II:900
resolves cleanly.

\begin{decisionbox}[Decision 3: $a$ is reserved for vortex core radius]
The symbol $a$ refers exclusively to the vortex core radius, which by
Paper~I's minimisation result is $a^* = \xi$ for the electron and
$a_p \sim \hbar/(m_p c)$ for the proton.  Where Paper~II currently uses
$a$ to mean a ring-scale or reconnection-scale quantity (II:900,
II:765--768), the symbol must be replaced by $R^*$ or by an explicit
ring-scale notation.
\end{decisionbox}

This implies:
\begin{itemize}
\item II:765 ($E_{\rm recon} \sim P_Q a^3$): $a$ here is the reconnection
      cavity scale, which is the ring scale $R^*$.  Should read
      $E_{\rm recon} \sim P_Q (R^*)^3$.
\item II:768 (``$a \sim \hbar_{\rm eff}/(m_e c)$''): if this is core radius,
      it equals $\xi$ (which equals $\hbar/(m_e c)$ at the electron
      reference mass), so the relation is correct but the use of
      $\hbar_{\rm eff}$ rather than $\hbar$ deserves attention.  See
      below.
\item II:900 (``$a \sim \xi/\alpha$''): this is the ring scale $R^*$,
      not the core radius.  Should read $R^* \sim \xi/\alpha$.
\end{itemize}

\subsection{Resolving $R^*$ and $\bar\lambda_e$}

These are unambiguous in the inventory.

\begin{decisionbox}[Decision 4: $R^*$ and $\bar\lambda_e$ stand]
\begin{align}
    R^*       &\equiv \xi/\alpha = \hbar/(\alpha m_e c)
                \approx 5.29\times 10^{-11}\,\mathrm{m}
                \tag{electron equilibrium ring radius, $\sim$ Bohr radius}\\
    \bar\lambda_e &\equiv \hbar/(m_e c) = \xi
                \tag{reduced electron Compton wavelength}
\end{align}
Note that $R^*$ is \emph{larger} than $\bar\lambda_e$ by a factor
$\alpha^{-1} \approx 137$.  The electron's spatial extent (its ring radius)
is therefore the Bohr radius scale, not the Compton scale.  This is
already implicit in Paper~I but worth stating since the relative size of
$R^*$ and $\bar\lambda_e$ matters for the g--2 form-factor argument.
\end{decisionbox}

\subsection{Resolving $m_{\rm grain}$: the heaviest stable defect}

The split in Decision~2 raises the question: what fixes
$m_{\rm grain}$?  Three candidates were considered in an earlier draft
of this audit:
\begin{enumerate}
\item $m_{\rm grain} \sim m_{\rm Planck}$ --- the SSV cutoff is the
      Planck scale, recovering the original $\xi \sim \ell_P$
      intuition that motivated the (now-removed) Paper~II line~415.
\item $m_{\rm grain} \sim m_H$ --- the Higgs amplitude-mode mass sets
      the medium grain.
\item $m_{\rm grain} = m_p$ --- the heaviest stable defect mass sets
      the medium's domain of validity.
\end{enumerate}

Candidate~(2) was the working position of an earlier audit draft.
It collapses on closer inspection because the Higgs is not a particle
but an event (see \S\ref{sec:higgs_event} below), and event energies
do not directly define structural length scales of the medium.

Candidate~(1) requires positing a Planck-scale microscopic
substrate, which SSV does not need to commit to (see
\S\ref{sec:domain}).

We therefore adopt candidate~(3):

\begin{decisionbox}[Decision 5: $m_{\rm grain}$ is the heaviest stable
                    defect mass]
\begin{equation*}
    m_{\rm grain} \equiv m_p \approx 1.673\times 10^{-27}\,\mathrm{kg}
    \;\;\Longrightarrow\;\;
    \ell_{\rm grain} \approx 2.10\times 10^{-16}\,\mathrm{m}.
\end{equation*}
The medium's domain of validity is the scale of its smallest stable
structure.  Below $\ell_{\rm grain}$, the LogSE description supports only
\emph{transient events}: denting events such as the $W$ boson, the $Z$
boson, the Higgs, the top quark, and the various heavy meson and baryon
resonances.  No new \emph{stable} defects exist below this scale.

This identification is empirical (the proton is observed to be the
heaviest stable defect) and provisional: if a heavier stable structure
were ever discovered, $m_{\rm grain}$ would update.  The framework's
mathematical machinery does not depend on the specific value of
$m_{\rm grain}$ above the proton scale; what matters is that
$m_{\rm grain}$ exists and finite.
\end{decisionbox}

This decision has the cleanest interpretation when combined with the
two other defect-spectrum facts:
\begin{itemize}
\item The lightest stable charged defect is the electron, with mass
      $m_e \approx 0.511\,\mathrm{MeV}/c^2$ (toroidal vortex ring).
\item The heaviest stable defect is the proton, with mass
      $m_p \approx 938\,\mathrm{MeV}/c^2$ (knotted breather).
\item The ratio $m_p/m_e \approx 1836$ is therefore the framework's
      defect spectrum width: the medium supports stable structures over
      three decades of mass, with no stable structures above
      $m_p$ or below $m_e$.
\end{itemize}

The hierarchy puzzle in SSV is then sharply posed: \emph{why does
the toroidal vortex ring have a major radius
$R^* = \xi/\alpha \approx 137\,\xi$ that is much larger than the
proton's knotted breather scale $a_p = \xi/1836$?}  This is now a
defect-spectrum question, not a defect-vs-canvas question.  The 3D
$\mathcal{L}+\mathcal{L}_\perp$ minimisation that determines the W
mass, the lepton generations, and the CP observables is also the
calculation that should fix the $m_e/m_p$ ratio.

\subsection{The Higgs as denting event, not as a particle or length scale}
\label{sec:higgs_event}

The audit's earlier draft identified $m_{\rm grain}$ with $m_H$ on the
grounds that the Higgs is the medium's amplitude mode.  This conflates
two distinct senses in which ``Higgs'' can refer to a feature of the
medium.

\paragraph{Two readings of ``the Higgs'' in SSV.}
\begin{itemize}
    \item \textbf{The amplitude mode.}  The medium has a longitudinal
          (amplitude) channel $\delta\rho$ in addition to the
          transverse phase channel $\delta\theta$.  This channel is
          gapped at the Higgs mass scale.  The amplitude mode is not a
          particle in the usual sense; it is a degree of freedom of the
          continuum medium, much as longitudinal sound is a degree of
          freedom of an ordinary fluid.  This sense is structural: it
          is a property of the LogSE Lagrangian, present at every point
          in the medium.
    \item \textbf{The 125\,GeV ``Higgs particle''.}  This is an
          \emph{event}: a localised compression of the order parameter
          $|\Psi|$ deep enough to ring the amplitude mode briefly
          before relaxing.  The number 125\,GeV is the energy required
          to drive a dent of a particular shape and depth ($|\Psi|\to 0$
          over a volume $V^*$) against the saturation pressure $P_0$,
          giving $E_H = P_0\,V^*$.  The geometry of the dent is set by
          the same 3D minimisation that determines the $W$ mass (where
          the analogous dent has a different cap geometry $V_{\rm cap}$).
\end{itemize}

The two readings have different status:
\begin{itemize}
\item The amplitude mode is a structural feature of the medium that
      exists everywhere; its mass scale (the gap) is a property of the
      Lagrangian and does \emph{not} fix a length scale of the
      medium's substrate.
\item The 125\,GeV event is a derived quantity --- a denting event
      energy, calculable in principle from $\{P_0, \alpha,$ medium
      geometry$\}$, with no separate fundamental status.
\end{itemize}

Neither reading licenses identifying $m_H$ with $m_{\rm grain}$.  The
amplitude-mode gap is a coupling-constant statement (which parameter of
the Lagrangian sets the gap), not a length-scale statement.  The
event-energy reading is a derived calculation, not a primitive of the
framework.

\paragraph{Implication for Paper~II \S5.3.}
The current text identifies the Higgs with the amplitude mode and
treats $\xi_{\rm EW} = \hbar/(m_H c)$ as a structural length of the
medium.  This needs revision.  $\xi_{\rm EW}$ is the de~Broglie
wavelength of the amplitude-mode quantum, not a structural feature of
the medium beyond that.  The ``canvas grain'' language of \S5.3 should
be retained but tied to $\ell_{\rm grain} \sim a_p$, not to
$\xi_{\rm EW}$.  The Higgs section should make clear that:
\begin{enumerate}
\item The amplitude mode is a degree of freedom of the medium with a
      gap at $m_H c^2 \approx 125\,\mathrm{GeV}$.
\item The 125\,GeV resonance observed at the LHC is a denting event ---
      the amplitude-mode analogue of the $W$ Rayleigh-jet event ---
      whose energy follows from $P_0$ and the dent geometry.
\item Neither of these identifies a structural length scale of the
      medium below $\ell_{\rm grain} \sim a_p$.
\item The Yukawa coupling $g_f = m_f/v$ retains its mechanical
      reading (defects displace the canvas; the coupling measures how
      much), without requiring a sub-grain Higgs length.
\end{enumerate}

\subsection{The framework's domain of validity}
\label{sec:domain}

The decisions above are clarified by an explicit statement of SSV's
domain of validity as an effective continuum theory.

\paragraph{The continuum analogy.}
The LogSE description of the saturated superfluid vacuum is a
\emph{continuum effective theory}, in the same epistemological category
as Navier--Stokes, the Madelung equations, or Landau theory of phase
transitions.  Such theories have three features:
\begin{enumerate}
\item They are predictive within their domain of validity.
\item They have a definite scale below which the description fails.
\item They are silent about what underlies them, and that silence is a
      feature, not a defect.
\end{enumerate}
Navier--Stokes describes water flow with extraordinary precision down
to scales of order the molecular spacing of water ($\sim 0.3\,$nm).
Below that scale, the equations fail; what is ``there'' is described
by atomic and quantum-mechanical theories that have no Navier--Stokes
character.  The transition from one description to the other is not a
boundary in space but a boundary in scale: at scales above
$\sim 1\,$nm, water is a fluid; at scales below, it is a system of
H$_2$O molecules; and at scales below that, atoms; and so on.  None of
these descriptions is wrong --- each has its domain.

\paragraph{The SSV analogue.}
SSV is the LogSE description of the saturated vacuum.  Its domain of
validity extends down to $\ell_{\rm grain} \sim a_p \approx 2\times
10^{-16}\,\mathrm{m}$, the scale of the smallest stable structure the
medium supports.  Within this domain, the framework makes sharp
predictions: the four forces, the Standard Model particle spectrum,
the discrete symmetries, gravitational and cosmological structure.

Below $\ell_{\rm grain}$, three things are true simultaneously:
\begin{itemize}
\item The medium continues to support \emph{transient events} ---
      denting events at the Higgs and $W/Z$ scales, top-quark events,
      heavy quarkonium states, and so on.  These are described by the
      same LogSE machinery applied to localised, non-stable
      configurations.
\item No new \emph{stable defects} exist.  The proton is the heaviest
      stable knotted breather; below its mass scale, all structures
      are transient.
\item The continuum description itself approaches the limit of its
      validity.  At sufficiently fine scales the question ``what is
      the medium made of?'' becomes well-posed, but the answer is not
      within SSV.
\end{itemize}

\paragraph{What is below the grain?}
The framework is deliberately silent.  Several candidate microscopic
substrates exist in the literature: string-theoretic worldsheets,
loop-quantum-gravity spin networks, causal-set elements, twistor-space
configurations, and others.  Any of these could in principle be the
microscopic theory whose continuum limit reproduces the LogSE.  SSV
does not commit to any of them.  This is a structural strength: the
framework's predictions in its domain of validity are independent of
which microscopic theory eventually proves correct, and conversely, any
viable microscopic theory must reproduce SSV's predictions in the
appropriate continuum limit.

\paragraph{Implication for the UV-cutoff argument (Paper~II \S6.6).}
The graviton-free / UV-divergence-dissolved argument now reads cleanly:
perturbative quantum gravity diverges because it treats the metric as
a fundamental field with arbitrary momentum modes.  In SSV, gravity is
a derived effect of medium dynamics, and medium dynamics is a continuum
description with a definite domain of validity.  Modes with
$k \gg 1/\ell_{\rm grain}$ are not regulated by hand --- they are
\emph{outside the framework's domain}, and assertions about them are
ill-posed within SSV in the same way assertions about turbulent
eddies smaller than a water molecule are ill-posed within
Navier--Stokes.  The divergent integrals of perturbative quantum
gravity therefore never appear: not because the framework supplies a
small-distance cutoff, but because the modes that would generate the
divergences are not part of the theory.

\subsection{Resolving the g--2 finite-ring correction (content issue)}
\label{sec:g2}

With $R^* \approx 137\,\bar\lambda_e$, the form factor
$F(kR^*) = J_0(kR^*)$ does \emph{not} admit the small-argument expansion
$1 - (kR^*)^2/4$ at the Compton scale $k \sim 1/\bar\lambda_e$.  At that
scale, $kR^* \approx 137$ and $J_0(137) \approx 0.07$ (oscillatory, not
small).

The correct statement is one of two things, and the choice has to be made
based on the physics:

\begin{enumerate}
\item \textbf{The form factor cuts the loop integral off at $k \sim 1/R^*$
      rather than at $k \sim 1/\bar\lambda_e$.}  In that case the loop
      integral is dominated by momenta well below the Compton scale, and
      the leading correction to $g - 2$ is suppressed by
      $(\bar\lambda_e/R^*)^2 = \alpha^2 \approx 5.3\times 10^{-5}$ relative
      to the leading $\alpha/\pi$ term.  Numerically:
      \[
        \delta(g-2)/2 \sim \alpha^3/\pi \sim 4\times 10^{-7},
      \]
      below current $(g-2)/2$ sensitivity of $\sim 10^{-13}$ by 6 orders
      of magnitude.  This is consistent with experiment.
\item \textbf{The relevant size scale for the loop is the core radius
      $a^* = \xi$, not the ring radius $R^*$.}  In that case the
      expansion variable is $(\xi/\bar\lambda_e)^2 = 1$ at the electron
      reference mass, and the correction is order $\alpha$ --- comparable
      to the leading Schwinger term.  This contradicts the
      observed agreement with QED at high precision and is therefore
      ruled out.
\end{enumerate}

We adopt resolution~(1) provisionally; a rigorous derivation of the loop
integral with the toroidal-vortex form factor is required to make this
firm.  This is flagged as an open computation.

\subsection{Resolving the UV-cutoff argument (content issue)}

The original text in Paper~II \S6.6 argues that the ultraviolet
divergence of perturbative quantum gravity dissolves because momentum
modes $k \gg 1/\xi$ do not exist in the medium.  With Decision~2,
$1/\xi \approx m_e c/\hbar$ is the electron Compton momentum --- a scale
well-explored in QED with no breakdown.  The argument as written is
therefore wrong.

The correct statement, developed in \S\ref{sec:domain} above, is that
modes with $k \gg 1/\ell_{\rm grain}$ are \emph{outside the framework's
domain of validity} rather than absent in some sharper sense.  With
Decision~5, $\ell_{\rm grain} \approx 2\times 10^{-16}\,\mathrm{m}$ is
the proton scale.  The graviton-free / UV-divergence-dissolved argument
becomes the same kind of argument as Navier--Stokes's silence about
sub-molecular eddies: the divergent integrals of perturbative quantum
gravity simply do not arise within SSV, because the modes that would
generate them are not part of the theory.

The argument is sharper than the original framing because it does not
require a Planck-scale cutoff or any specific microscopic substrate.
It only requires that SSV be an effective continuum theory with a
domain of validity, which is automatic for any continuum theory.

\subsection{Resolving $\xi_{\rm EW}$}

With Decisions~2 and 5 in place, $\xi_{\rm EW} \equiv \hbar/(m_H c)
\approx 1.58\times 10^{-18}\,\mathrm{m}$ is the de~Broglie wavelength of
the amplitude-mode quantum.  It is \emph{below} $\ell_{\rm grain}$, in
the regime where the medium supports only transient events.  This is
consistent with the Higgs being a denting event (\S\ref{sec:higgs_event})
rather than a structural length of the medium: $\xi_{\rm EW}$ is the
characteristic scale of the dent, not a substrate scale.

The numerical hierarchy $\xi_{\rm EW} \ll \ell_{\rm grain} \ll \xi$
(roughly $10^{-18}\,\mathrm{m} \ll 10^{-16}\,\mathrm{m} \ll
10^{-13}\,\mathrm{m}$) is consistent: the Higgs event is a deeper
dent than the proton's stable knot, which is in turn smaller than the
electron's stable ring.  All three numbers fall out of the same
$\mathcal{L}+\mathcal{L}_\perp$ machinery, applied to different
configurations.

% ═══════════════════════════════════════════════════════════════════════
\section{Defined-Symbol Reference Table}
% ═══════════════════════════════════════════════════════════════════════

The following table gives the canonical definitions to be used in revised
versions of Papers~I and II.  Numerical values use CODATA fundamental
constants and the framework's identification $m_0 = m_e$ for the electron
defect (Paper~I, line~322--323).

\renewcommand{\arraystretch}{1.6}
\begin{longtable}{@{}p{2.5cm}p{4.5cm}p{3.5cm}p{4cm}@{}}
\toprule
\textbf{Symbol} & \textbf{Definition} & \textbf{Numerical value}
                & \textbf{First use / role} \\
\midrule
\endhead
$m_e$ & Electron rest mass & $9.109\times 10^{-31}$\,kg & 
        Empirical input (Paper~I) \\
$m_p$ & Proton rest mass & $1.673\times 10^{-27}$\,kg & 
        Derived from Paper~I ladder \\
$m_H$ & Higgs amplitude-mode gap (denting-event energy) &
        $\approx 125\,\mathrm{GeV}/c^2$ &
        Empirical (Paper~II \S5.3); not a length scale of the medium \\
$m_{\rm grain}$ & Heaviest stable defect mass; sets domain of validity &
        $= m_p \approx 938\,\mathrm{MeV}/c^2$ &
        Decision~5 of this audit; UV cutoff (Paper~II \S6.6, revised) \\
$m_0$ & Reference mass for a given defect's healing-length formula
        & $m_e$ for electron, $m_p$ for proton, etc. &
        Defect-relative (Paper~I line~322--323) \\
\midrule
$\alpha$ & Fine-structure constant; $\alpha = c_\perp^2/c^2$ & $1/137.036$ &
        Squared ratio of transverse to longitudinal wave speeds (Paper~I) \\
$c$ & Longitudinal sound speed $=$ light speed &
        $2.998\times 10^8$\,m/s & Stiffness fixing (Paper~I \S2) \\
$c_\perp$ & Transverse wave speed $= \sqrt\alpha\,c$ &
        $2.561\times 10^7$\,m/s & Photon dispersion (Paper~II \S3.3) \\
\midrule
$\xi$ & Electron healing length $\equiv \hbar/(m_e c)$ &
        $3.862\times 10^{-13}$\,m & 
        Vortex core size for electrons (Paper~I, recurring) \\
$\ell_{\rm grain}$ & Medium domain-of-validity scale $\equiv \hbar/(m_p c)$ &
        $\approx 2.10\times 10^{-16}\,\mathrm{m}$ & 
        Smallest scale at which the LogSE description is valid; below
        this, only transient events \\
$\bar\lambda_e$ & Reduced electron Compton wavelength $\equiv \hbar/(m_ec)$
                  $= \xi$ &
        $3.862\times 10^{-13}$\,m & QED reference scale (Paper~II \S3.5) \\
$a$ & Vortex core radius (generic) &
        $a^* = \xi$ for electron;
        $a_p = \hbar/(m_p c) \approx 2.10\times 10^{-16}$\,m for proton &
        Reserved for core only; not for ring or reconnection scales
        (Paper~I \S3) \\
$R^*$ & Electron equilibrium ring radius $\equiv \xi/\alpha = \hbar/(\alpha m_ec)$ &
        $5.292\times 10^{-11}$\,m & Bohr radius scale (Paper~I, boxed) \\
$\xi_{\rm EW}$ & Higgs amplitude-mode wavelength $\equiv \hbar/(m_H c)$ &
        $1.578\times 10^{-18}$\,m & Characteristic scale of the Higgs
        denting event; below $\ell_{\rm grain}$, in the
        transient-event regime (Paper~II \S5.3, revised) \\
$\ell_P$ & Standard Planck length $\equiv \sqrt{\hbar G/c^3}$ &
        $1.616\times 10^{-35}$\,m &
        Reference comparison only; SSV is silent about scales below
        $\ell_{\rm grain}$ and does not commit to Planck-scale
        microphysics (\S\ref{sec:domain}) \\
\midrule
$\mu_0$ & Base mass-ladder rung $\equiv m_e/\alpha$ &
        $\approx 70.025$\,MeV/$c^2$ &
        Pion/lepton ladder (Paper~I) \\
$\kappa_0$ & Circulation quantum $\equiv h/m_0$ &
        $h/m_e$ for electron defect &
        Quantised circulation (Paper~I) \\
\midrule
$\rho_0$ & Saturated background density &
        Determined by $b = m_0 c^2/(2\rho_0)$ &
        Medium parameter (Paper~I) \\
$P_0$ & Saturation pressure $\approx \rho_0 c^2$ &
        Determined by $\rho_0$ &
        W cap energy (Paper~II \S5.4); Yukawa weighting (\S5.3) \\
$\lambda$ & Chiral-shear coupling $\equiv \alpha^2 \rho_0/m_0$ &
        Per-defect via $m_0$ &
        Sets $c_\perp$ (Paper~I) \\
\bottomrule
\end{longtable}

\subsection{Things this table makes explicit}

\begin{enumerate}
\item $\xi$ and $\bar\lambda_e$ are numerically the same quantity at the
      electron reference mass; using both symbols is redundant and a source
      of confusion.  Recommendation: keep $\xi$ for hydrodynamic contexts
      (vortex core, healing length, sub-grain coupling), keep
      $\bar\lambda_e$ for QED-like contexts (form-factor expansion variable,
      g--2 derivation).  Note in the paper that they coincide.
\item $a^* = \xi$ for the electron defect (Paper~I minimisation).  When
      Paper~II uses $a$ at the ring scale ($\sim \xi/\alpha$), it should
      use $R^*$ instead.
\item $\ell_{\rm grain} \approx 2\times 10^{-16}\,\mathrm{m}$ is the
      proton-scale domain-of-validity boundary, introduced by Decision~5.
      It should appear explicitly in Paper~II \S6.6 in place of the
      current misuse of $\xi$.  The hierarchy
      $\xi_{\rm EW} \ll \ell_{\rm grain} \ll \xi$ separates transient
      events ($\xi_{\rm EW}$ regime) from stable defects (between
      $\ell_{\rm grain}$ and $\xi$).
\item $m_0$ is defect-relative.  This is already in Paper~I, but the
      consequences for Paper~II are not currently spelled out.
\item $m_H$ is not a length scale of the medium.  $\xi_{\rm EW}$ is the
      characteristic scale of the Higgs denting event, not a substrate
      property.  The Higgs is a feature of the medium's amplitude
      channel, expressed as both a coupling-constant statement (the gap)
      and a derived event-energy (the LHC resonance), but neither
      reading licenses identifying $m_H$ with $m_{\rm grain}$.
\end{enumerate}

% ═══════════════════════════════════════════════════════════════════════
\section{What Counts as a Propagation Speed}
\label{sec:propagation_scope}
% ═══════════════════════════════════════════════════════════════════════

The audit's discussion of $c$ and $c_\perp$ concerns \emph{propagation
speeds} of localised disturbances in the medium: a transverse phase
wiggle moves at $c_\perp = \sqrt\alpha\,c$, a longitudinal density
wiggle moves at $c$.  These are mechanical wave speeds in the same sense
that sound speed in water is a mechanical wave speed.

Not every physical phenomenon associated with the medium falls into this
category.  In particular, correlations between spatially separated
parts of the medium's saturated state are properties of the
\emph{global state} rather than of any localised disturbance, and do
not have a propagation speed at all --- they are not signals.  This
includes the correlations responsible for quantum entanglement: the
saturated medium is a single coherent state, and what experiment
detects as ``spooky action at a distance'' is the structural fact that
this state is whole, not the propagation of any influence between its
parts.

\begin{notebox}[Acknowledgement, not a programme]
SSV's saturated medium is a coherent global object.  Correlations within
that global object --- including those exhibited by entangled particles
--- have a holographic character: they are properties of the medium's
whole-state representation rather than of localised disturbances
travelling through it.  The lower bounds on ``the speed of spooky action
at a distance'' obtained by Salart \emph{et al.}\ (2008, $\geq 10^4 c$)
and Yin \emph{et al.}\ (2013, $\geq 10^4 c$ with locality loopholes
closed) constrain hypothetical signal velocities; they are vacuous in
the present framework because no signal is postulated.

The audit acknowledges this aspect of the medium's structure without
developing it.  A full treatment of the holographic content of the
saturated state --- the relationship between SSV and the holographic
principle, how entanglement entropy maps onto medium properties, and
whether positive predictions can be extracted (e.g.\ entanglement
behaviour in strong gravitational fields, where $\rho$ deviates
significantly from $\rho_0$) --- is left for future work, possibly as
a separate paper in the series.

The single point this audit insists on is that propagation-speed
questions are distinct from holographic-correlation questions.
Conflating the two has caused confusion in earlier drafts; keeping
them separate clarifies what the framework's speed claims are
actually about.
\end{notebox}

The Paper~II \S5.6 (\emph{CP and the Strong CP Problem in a Medium That
Remembers}) discussion of CP and the medium's recorded
history is structurally adjacent: both treatments rest on the medium
being a single coherent whole rather than a collection of local
properties.  A one-sentence forward-pointer in that subsection
acknowledging the holographic aspect is appropriate; the present
audit's stance is that no further development is needed in either
Paper~I or Paper~II.

\subsection{The radiation channel between defects}
\label{sec:radiation_channel}

The discussion above contrasted two \emph{propagating} medium modes
(transverse at $c_\perp$, longitudinal at $c$).  A separate question is
which of these modes carries radiation \emph{between defects}.  In
SSV, defects are the only sources of observable signals: photons,
gravitons (in the standard sense), and any other ``messenger field''
of conventional physics are either identified with medium modes or do
not exist as separate objects.  The framework's strongest claim about
the latter is that there is \emph{no graviton}: gravity is not a
distinct field with its own propagation speed but an interaction
between defects mediated by their pressure responses in the medium.

This recasts the radiation-channel question.  In a featureless
linear medium, both transverse and longitudinal modes can in principle
carry energy outward.  In SSV the relevant radiation is sourced by
\emph{defects} --- localised, mass-energy-carrying topological
structures --- and the channel that carries their outward signal is the
one to which their multipole moments couple.  Two structural facts
fix the answer:

\begin{enumerate}
\item \textbf{Mass produces a density depression, not a compression.}
      A defect's stable presence rarefies the medium locally
      ($\rho < \rho_0$ near a mass; Paper~II \S6).  A time-varying
      depression pattern is not a longitudinal compression wave; it is
      a transverse deformation of the saturated state propagating
      outward.
\item \textbf{The longitudinal channel is gapped.}  The longitudinal
      compression mode carries the Higgs amplitude excitation, with
      gap $m_H c^2 \approx 125$\,GeV (Paper~II \S5.3).  Sources
      radiating at any frequency $\hbar\omega \ll m_H c^2$ ---
      including all astrophysical sources of gravitational radiation
      ever observed --- cannot excite this gapped mode.  The
      longitudinal channel is energetically inaccessible to ordinary
      defect-defect interaction.
\end{enumerate}

The combined consequence is sharp: \textbf{all defect-defect radiation
in SSV propagates in the transverse channel, at $c_\perp$.}  This
includes electromagnetic radiation (the chirality-coupled response,
Paper~II \S3) and gravitational radiation (the mass-energy-coupled
response, Paper~II \S6).  The two are different multipole patterns
in the same channel, not different channels.

\begin{decisionbox}[Decision 6: Single radiation channel for defect-defect signals]
All radiation between defects in SSV --- electromagnetic, gravitational,
and any other observable signal carried between localised mass-energy
configurations --- propagates at $c_\perp = \sqrt{\alpha}\,c$.  The
longitudinal mode at $c$ exists as a structural property of the
medium (it sets the saturation pressure $P_0$, the Higgs amplitude
gap, and the near-field dent geometry of $W$, $Z$, top and Higgs
events), but it is gapped at $m_H c^2$ and is not excited by ordinary
astrophysical sources.  In particular, gravitational waves and
electromagnetic radiation from astrophysical sources arrive at any
detector simultaneously up to source-physics timescales, with no
propagation-speed difference.

This resolves the apparent GW170817 problem.  The framework's
$\alpha = c_\perp^2/c^2$ identification is preserved; what changes is
the recognition that $c$ is an internal medium scale not directly
observed as a propagation speed of any astrophysical signal.
\end{decisionbox}

% ═══════════════════════════════════════════════════════════════════════
\section{Open Issues}
% ═══════════════════════════════════════════════════════════════════════

The following are items the audit identified that are \emph{not} notation
issues but genuine physics questions, and which the canonical table cannot
itself answer.

\begin{enumerate}
\item \textbf{Targeted revision of Paper~II \S6 (Bjerknes mechanism).}
      The current \S6 derivation treats gravitational radiation as
      longitudinal acoustic radiation between pulsating breathers.
      Decision~6 of \S\ref{sec:radiation_channel} requires this be
      reframed: the time-averaged $1/r^2$ Bjerknes force survives, but
      the time-varying radiation channel is transverse rather than
      longitudinal.  This is a notational and interpretive update to
      the existing derivation, not a structural change --- the
      mechanism remains ``pressure-response between pulsating
      sources,'' with the carrier mode now correctly identified.

      The motivation is empirical: GW170817 (binary neutron-star
      merger, 2017) showed gravitational and electromagnetic signals
      arriving within $\sim$1.7\,s over $\sim$130\,Mly, constraining
      $|\Delta v|/v \lesssim 10^{-15}$.  The mainstream astrophysical
      interpretation attributes the 1.7\,s entirely to source physics
      (relativistic-jet launch from the merger remnant, jet breakout
      through the surrounding ejecta, and propagation to the
      photospheric radius); the propagation-speed difference is
      consistent with zero.  Decision~6 is consistent with this
      observation: gravity and light share the transverse channel and
      arrive simultaneously up to source-physics timescales.  No
      further empirical accommodation is required, but the \S6 text
      must be brought into line with the radiation-channel
      identification.

\item \textbf{Derivation of $m_p$ (and hence $m_{\rm grain}$) from
      $\{m_e, \alpha\}$.}  Paper~I states only $m_e$ and $\alpha$ as
      empirical inputs.  If this claim is to be sustained, $m_p$ must be
      derivable from the framework rather than empirical.  The current
      Paper~I uses $m_p = N_p\,m_e/\alpha$ with $N_p \approx 13.44$
      determined empirically, with the form factor $F$ expected from a
      3D minimisation.  Until that minimisation is performed, the claim
      that the framework has only two empirical inputs is provisional:
      $m_p$ functions as a third empirical input via $N_p$.
\item \textbf{The g--2 finite-ring correction (\S\ref{sec:g2}).}
      Resolution~(1) gives a parametrically suppressed correction,
      but the loop integral with the toroidal-vortex form factor has
      not been done rigorously.  Until it is, the $\sim 10^{-44}$
      estimate currently in Paper~II is not defensible at any level.
\item \textbf{$\hbar_{\rm eff}$ vs.\ $\hbar$.}  Paper~II line~768 uses
      $\hbar_{\rm eff}$ in the W-mass scalar estimate.  It is not clear
      whether this is meant to be a separate symbol (a renormalised
      $\hbar$ in the medium) or a redundant notation.  This was not on
      the original audit list but came up; deserves clarification.
\item \textbf{The proton's $m_0$.}  Decision~1 says the proton breather
      uses $m_0 = m_p$.  This means the proton's own healing length is
      $\hbar/(m_p c) = a_p$, which equals $\ell_{\rm grain}$.  This is
      not a coincidence: the proton's reference mass is the medium's
      domain-of-validity scale precisely because the proton is the
      heaviest stable defect.  When gravity sub-grain coupling is
      computed via $a_p/\xi = m_e/m_p$, the $\xi$ in the denominator is
      the \emph{electron} healing length (i.e.\ the electron's $m_0$
      value), not the proton's.  This cross-defect comparison should be
      flagged in Paper~II's gravity section as a deliberate use of two
      reference masses.
\item \textbf{The defect-spectrum hierarchy.}  Why does the toroidal
      vortex ring (electron) have $R^* \approx 137\,\xi$ while the
      knotted breather (proton) has $a_p \approx \xi/1836$?  Both
      arise from the same LogSE+chiral functional applied to different
      topologies.  The 3D minimisation that produces the W mass and
      the lepton generations should also produce this ratio.  This is
      the framework's hierarchy puzzle, sharply posed.
\item \textbf{The microscopic substrate.}  SSV is a continuum effective
      theory; its silence about ``what's below $\ell_{\rm grain}$'' is
      a feature, not a defect (\S\ref{sec:domain}).  Multiple candidate
      microscopic theories exist (string-theoretic worldsheets,
      loop-quantum-gravity spin networks, causal sets, twistor-space
      configurations).  SSV does not commit; any viable microscopic
      theory must reproduce SSV's continuum predictions.  This is
      flagged as an explicit \emph{non-deliverable} of the SSV
      programme: the framework neither needs nor provides a
      microscopic foundation.
\end{enumerate}

% ═══════════════════════════════════════════════════════════════════════
\section{Recommended Order of Manuscript Revisions}
% ═══════════════════════════════════════════════════════════════════════

Once Decisions~1--5 are accepted, the manuscript revisions are:

\begin{enumerate}
\item \textbf{Paper~I, line~322--323.}  Strengthen the ``$m_0$ is a
      reference mass scale (self-consistently $m_e$ for the electron)''
      declaration into a full paragraph.  State explicitly that $m_0$ is
      defect-relative, cite the proton case as the second example, and
      forward-pointer to Paper~II's use of $m_0 = m_p$ for the proton
      breather.  This is the single most important upgrade because
      every downstream issue stems from the under-specification of $m_0$.
\item \textbf{Paper~II \S2 (Framework recap).}  Add the symbol table
      (or a compact version of it) to the recap.  Distinguish $\xi$
      (electron healing length) from $\ell_{\rm grain}$ (medium
      domain-of-validity scale, identified with the proton scale by
      Decision~5).  Add a brief paragraph on the framework as a
      continuum effective theory with a definite domain of validity.
\item \textbf{Paper~II \S3.5 (g--2 finite ring correction).}  Replace
      the $\xi \to \ell_P$ substitution.  Present the correction in
      whichever resolution (1)--(2) of \S\ref{sec:g2} above is adopted,
      with a gapbox flagging that the rigorous loop integral is open.
\item \textbf{Paper~II \S5.3 (Higgs as canvas).}  Substantial revision.
      Distinguish:
      \begin{itemize}
      \item The amplitude mode as a structural feature of the medium
            (a property of the LogSE Lagrangian, with a gap at $m_H c^2$).
      \item The 125\,GeV LHC resonance as a denting event ($E_H = P_0
            V^*$ with $V^*$ from the same 3D minimisation that produces
            the $W$ mass).
      \end{itemize}
      Drop the framing that identifies $\xi_{\rm EW}$ as a structural
      ``canvas grain.''  Reframe the hierarchy puzzle as the defect
      spectrum question: why $m_e \ll m_p$, geometrically
      $R^* \approx 137\,\xi$ vs.\ $a_p \approx \xi/1836$.
\item \textbf{Paper~II \S5.4 (Rayleigh jet) and \S5.2 (W mass scalar).}
      Replace ring-scale uses of $a$ with $R^*$.  Reserve $a$ for core.
\item \textbf{Paper~II \S6.6 (graviton-free / UV cutoff).}
      Substantial revision.  Replace $\xi$ with $\ell_{\rm grain}$ in
      the cutoff argument.  Reframe the argument: divergent integrals
      of perturbative quantum gravity do not arise in SSV not because
      of a Planck-scale cutoff, but because the modes that would
      generate them are outside the framework's domain of validity ---
      the same kind of statement Navier--Stokes makes about
      sub-molecular eddies.  Drop the implicit Planck-scale framing.
\item \textbf{Paper~II abstract and conclusion.}  Reflect the reframed
      Higgs treatment, the domain-of-validity scope, and the sharpened
      hierarchy puzzle.  In particular, the abstract should state that
      the framework is a continuum effective theory of the saturated
      vacuum, valid down to $\ell_{\rm grain} \sim a_p$, with the
      microscopic substrate beneath that scale a separate question
      handed off to a future programme.

\item \textbf{Paper~II \S5.6 (CP and the medium that records).}  Add
      a one-sentence forward-pointer noting that the same coherent
      whole-state property of the medium that breaks process-symmetries
      also accommodates the holographic correlations of entangled
      states, with the speed-of-entanglement question outside the
      propagation-speed framing of \S\S3 and 6 (see audit
      \S\ref{sec:propagation_scope}).  No further development.

\item \textbf{Paper~II \S6 (gravity) and abstract.}  Targeted update
      reframing gravitational radiation as a transverse-channel
      ($c_\perp$) signal rather than a longitudinal ($c$) one.  The
      Bjerknes derivation is preserved structurally --- the
      time-averaged $1/r^2$ force, the proton sub-grain coupling, the
      $G$ consistency check, the no-graviton claim, and the
      domain-of-validity argument all stand unchanged.  What changes
      is the radiation channel that carries the time-varying part of
      the breather's pressure response: it is identified as
      transverse, on the grounds that (i) mass produces a density
      depression (a transverse deformation of the saturated state),
      not a compression, and (ii) the longitudinal channel is gapped
      at the Higgs amplitude scale and inaccessible to
      sub-electroweak sources.

      A short paragraph should be added to \S6 making explicit that
      gravitational and electromagnetic radiation share the same
      carrier channel (the transverse phase channel of the medium's
      response to defects), differing only in their multipole-source
      structure (mass-energy quadrupole vs.\ chirality dipole).  This
      is consistent with GW170817; the observed $\sim$1.7\,s delay
      between gravitational and electromagnetic signals from the
      neutron-star merger is attributed to source physics
      (relativistic-jet launch and breakout) rather than to any
      propagation-speed difference, in agreement with mainstream
      astrophysical interpretation.

      The abstract should reflect that gravity and electromagnetism
      share the radiation channel; the longitudinal mode at $c$
      carries the Higgs amplitude excitation and is structural to the
      medium, but is not the carrier of any astrophysical signal.
\end{enumerate}

This audit document, once approved, becomes the canonical reference for
those revisions.

\end{document}
